\chapter{Background: Hardware Description, State Formalisms, and Verification}\label{ch02:background}

\section{Purpose and reading discipline of this survey}\label{ch02:purpose}

This chapter assembles, at tutorial depth, the established engineering and
mathematical background on which the rest of the monograph stands: hardware
description languages; explicit state formalisms; the industrial practice of
netlists and rule checking; graph algorithms for circuit analysis; the
distinction between simulation verification and operational validation;
property-based testing; provenance modeling; arithmetic control; and
information geometry on probability simplices. The survey is self-contained
because the monograph's readership spans communities --- cognitive-science
readers who have never read a netlist, software readers who have never
derived a Fisher metric, hardware readers who have never needed a formal
provenance vocabulary --- and the later chapters assume all of this
vocabulary at once.

Two disciplines govern the reading of every section below.

\begin{boundary}[Background is background]\label{ch02:bnd:background}
Nothing in this chapter is a contribution of the CAS PCBA 8 project. Every
concept, formalism, algorithm, standard, and theorem surveyed here is prior
art of its own community, cited where a citation exists in this monograph's
reference list and described in general terms where the material is
standard industrial practice. Conversely, nothing in this chapter validates
CAS: the fact that VHDL, Statecharts, design-rule checking, Tarjan's
algorithm, verification-and-validation discipline, QuickCheck, PROV, IEEE
754, or the Fisher--Rao metric are sound and useful in their home domains
is evidence about \emph{those} domains and about engineering method. It is
not evidence that the CAS ontology, its cognitive interpretation, or any
CAS PCBA 8 design choice is empirically correct. The support this
literature lends is at the level of method only.
\end{boundary}

The second discipline is a map. For each background item, the section
states (i) what the item is, (ii) what its home community learned from it,
and (iii) which mechanism of the CAS PCBA 8 architecture borrows the lesson
--- always as an organizational borrowing, never as a transferred validity
claim. Section~\ref{ch02:summary} collects the map in one table.

\section{Hardware description languages}\label{ch02:hdl}

\subsection{From annotated drawings to machine-readable design text}\label{ch02:hdl:drawings}

Before the mid-1980s, the dominant medium of digital-hardware design
description was the schematic drawing supplemented by prose. The medium had
the same epistemic weaknesses that Chapter~\ref{ch01:introduction} attributed
to informal cognitive diagrams: a drawing fixes topology but leaves timing,
interface contracts, allowed values, and intended behavior to convention and
caption; two engineering teams could read the same drawing into different
circuits; and there was no mechanical way to ask a drawing whether it was
even self-consistent. As designs grew from hundreds to millions of gates,
the medium failed to scale in three independent ways: description (the
drawing could not be complete), communication (the drawing could not be
unambiguous), and verification (the drawing could not be executed or
checked).

Hardware description languages (HDLs) answered all three failures at once by
making the design text itself formal, machine-readable, and executable. An
HDL description is not documentation \emph{about} a design; under the
language's semantics, it \emph{is} the design, and every downstream activity
--- simulation, synthesis, rule checking, equivalence checking, archiving
--- consumes the same text. This is the single most consequential idea the
present monograph borrows: the move from diagram-plus-prose to a typed
design language with defined semantics is a methodological act, independent
of subject matter, whose payoff is that design intent becomes checkable.

\subsection{What VHDL demonstrated}\label{ch02:hdl:vhdl}

VHDL, standardized by the IEEE as Std 1076 and maintained through successive
revisions to the current language reference manual \cite{vhdl}, is the
canonical demonstration of this move, and several of its specific design
decisions are instructive far beyond hardware.

First, VHDL separates \emph{interface} from \emph{implementation} as a
matter of language structure: an entity declaration fixes the ports ---
names, directions, and types --- through which a design unit may be used,
while one or more architecture bodies supply behavior. A component is usable
only through its declared interface, and the language makes it impossible to
connect to a port that was not declared. The lesson is interface
discipline: contracts live in the text, not in the reader's head.

Second, VHDL is strongly and expressively typed. Signals and variables
carry declared types; operations between incompatible types are rejected at
analysis time rather than misinterpreted at run time; and the language's
package mechanism lets a community standardize shared types with resolved
semantics. Decades of industrial experience taught that the type checker's
refusals are not bureaucratic friction but early defect detection: a large
class of wiring and width errors that would otherwise surface as silent
simulation misbehavior surface instead as analysis errors with source
locations. The lesson is that type safety is a defect-prevention
technology.

Third, VHDL has defined simulation semantics --- an event-driven model with
delta cycles, in which the meaning of a description is the trace of signal
transactions it produces under a defined scheduling discipline. Because the
semantics is part of the standard, two conforming simulators must agree on
what a description means, and a simulation result is a property of the
description rather than of one vendor's tool. The lesson is that
executability with defined semantics converts design text into a
replayable object.

Fourth, VHDL demonstrated the value of \emph{standardization of description
itself}. The language emerged from a documentation requirement --- the need
for a uniform, tool-independent way to describe hardware across
contractors and decades --- and its longevity shows that a stable,
versioned, publicly specified description language is an archival
technology: a design written down in it remains interpretable when the
tools, companies, and engineers of its era are gone.

Fifth --- a point easily missed from outside the field --- VHDL culture
developed the \emph{testbench} as a first-class design artifact: a
description, written in the same language, whose sole purpose is to
exercise another description and check its responses. The design under
test and the apparatus that tests it share a type system, a simulator, and
an archive, so a verification claim (``this entity passed this testbench'')
is a statement about two identified artifacts and can be rerun by anyone
holding both. The testbench idea generalizes far beyond hardware: it is
the recognition that checking apparatus deserves the same rigor of
description, identity, and versioning as the thing checked. CAS PCBA 8's
fixtures --- including, notably, its known-fault fixtures, which encode
designs that \emph{must} fail the checker in a declared way --- are
testbench culture applied to a rule checker.

Sixth, the industry's experience with the \emph{synthesis subset} carries
a cautionary lesson in the other direction. Not every legal VHDL
description can be synthesized into hardware; synthesizable VHDL is a
subset, and for years the boundary of that subset was tool-dependent and
imperfectly documented, so that a description could simulate one behavior
and synthesize another --- the simulation/synthesis mismatch, a notorious
class of defect. The lesson is that when a language is richer than the
machinery that consumes it, the gap between what can be \emph{written}
and what is actually \emph{supported} must be explicit, versioned, and
checked, or it becomes a source of silent divergence. This lesson bears
directly on the monograph's central adverse finding: the source
baseline's published grammar covers only a fraction of the constructs its
own reference designs use, an unmanaged written/supported gap of exactly
the kind synthesis history warns about. The candidate's response ---
declare the supported surface exactly, fail closed outside it, and demand
that any future grammar publish a construct-coverage matrix --- is the
hardware community's remedy, transplanted.

For this monograph, VHDL's demonstration is cited strictly at the level of
method \cite{vhdl}: formal machine-readable design description with typed
interfaces and defined execution semantics is feasible, scalable, and
industrially transformative. CAS PCBA 8's typed board model, closed
component catalog, and deterministic simulator borrow the organization.
The borrowing carries no validity: a cognitive-system hypothesis written in
a typed board model is exactly as conjectural as the same hypothesis on a
whiteboard --- it is merely explicit, checkable, and replayable in a way
the whiteboard version is not.

\subsection{Description versus realization}\label{ch02:hdl:realization}

One further HDL lesson deserves its own subsection because it prefigures
the monograph's first load-bearing separation. An HDL description that
simulates correctly is not thereby a working chip. Between description and
silicon lie synthesis, technology mapping, layout, parasitics, process
variation, and physical verification, and hardware culture maintains a
strict vocabulary boundary between what simulation shows (the description's
behavior under a model) and what characterization of fabricated parts shows
(the artifact's behavior in the world). The industry learned, sometimes
expensively, that eliding this boundary --- treating simulation sign-off as
silicon sign-off --- is a category error with recall-sized consequences.
The model-versus-world separation of
Section~\ref{ch01:sep:modelworld} is this hardware lesson, imported whole.

\section{Statecharts and explicit state semantics}\label{ch02:statecharts}

\subsection{The problem: state explosion and implicit dynamics}\label{ch02:statecharts:problem}

Finite-state machines are the classical formalism for reactive behavior,
but flat state machines fail on realistic systems for a structural reason:
the number of global states is the product of the states of the system's
independently varying parts, so a handful of concurrent features multiplies
into an unreadable and unmaintainable diagram. Practitioners of the 1970s
and early 1980s responded by leaving most of the state structure implicit
--- in prose, in code, in convention --- which reproduced, for dynamics,
exactly the auditability failure that schematics-plus-prose produced for
structure.

\subsection{Harel's contribution}\label{ch02:statecharts:harel}

Harel's Statecharts \cite{harel} attacked the problem by enriching the
visual formalism itself with three orthogonal mechanisms, each with defined
semantics. \emph{Hierarchy} (depth): states may contain substates, and a
transition from a superstate abstracts a bundle of transitions from all its
substates, so the diagram supports abstraction and refinement without
losing precision. \emph{Orthogonality} (concurrency): a state may be the
AND-composition of independent regions, so the diagram represents a product
of $n$ regions in size proportional to their sum, taming the exponential
blowup syntactically. \emph{Broadcast communication}: events raised in one
region are visible to others, giving a defined mechanism for coupling
concurrent parts. Harel's summary formula --- statecharts = state-diagrams
+ depth + orthogonality + broadcast-communication --- names precisely the
three additions.

The deeper lesson of the Statecharts program, however, is not any single
mechanism but the demand it institutionalized: \emph{a visual formalism
must have a semantics}. A diagram whose boxes and arrows lack defined
meaning is not a specification, however suggestive; once depth,
concurrency, and communication have defined semantics, the diagram becomes
a mathematical object that tools can execute, check, and transform. The
subsequent absorption of statechart variants into mainstream modeling
practice testifies that the demand was answerable at industrial scale.

\subsection{The aftermath: semantics is where the work is}\label{ch02:statecharts:aftermath}

The years after the original paper taught a second lesson that the first
had made possible. Once statecharts were executable, the community
discovered that the notation admitted many defensible but different
semantics: whether transitions taken in one step can trigger further
transitions in the same step; how conflicts between a superstate's
transition and a substate's are resolved; whether the effects of
simultaneous events are perceived instantaneously or in the next step;
what history connectors restore, exactly, on re-entry. Dozens of published
variants differed on such points, and two tools could disagree about the
behavior of the same chart while both plausibly claiming to implement
``statecharts.'' The lesson is not that the formalism failed --- it is
that formalization \emph{surfaced} ambiguities which informal boxes and
arrows had contained invisibly all along, and that the surfacing is the
value. A notation without semantics cannot even disagree with itself;
executable semantics converts hidden ambiguity into visible, decidable,
versionable choice. For the present work the corollary is direct: a
modeling language for cognitive hypotheses must expect that its early
versions under-specify, must version its semantic choices explicitly, and
must treat a discovered ambiguity (such as a grammar that does not say
what its own corpus means) as a first-class finding to record rather than
a nuance to gloss.

\subsection{The lesson borrowed}\label{ch02:statecharts:lesson}

CAS PCBA 8 borrows the demand, not the notation. Its board model requires
every component to declare statefulness and time semantics explicitly;
components without declared state semantics are unsupported by the
simulator and fail closed rather than defaulting to an implicit dynamics;
feedback --- the circuit-shaped analog of coupled concurrent regions ---
must be declared, delayed, and bounded rather than left emergent. Where the
architecture currently supports only a narrow deterministic profile
(stateless transfer with declared unit delays), it says so, and richer
state semantics are recorded as open obligations rather than assumed. The
Statecharts literature \cite{harel} is cited for the methodological point
that explicit state, hierarchy, concurrency, and communication semantics
are what turn behavioral diagrams into checkable objects; it neither
knows of nor validates any cognitive interpretation.

\section{Electronic design automation practice: netlists and rule checking}\label{ch02:eda}

The material in this section is standard industrial practice, described in
general terms; it is the accumulated craft of the electronic design
automation (EDA) industry rather than the contribution of any single
citable source.

\subsection{The netlist as central data structure}\label{ch02:eda:netlist}

At the center of every EDA flow sits the netlist: a complete, queryable
enumeration of a design's components (with their kinds and parameters),
their ports, and the nets --- the connection equivalence classes --- that
join port to port. The netlist is the design's connectivity reduced to a
data structure, and its centrality is a methodological fact: once
connectivity is a data structure rather than a drawing, connectivity
questions become computations. Which inputs are unconnected? Which nets
have two drivers? Which component instances share a reference designator?
What is the fan-out of this output? Every such question, undecidable
against a pile of schematic sheets by anything but human inspection, is a
traversal over the netlist.

The netlist also established an identity discipline. Each component
instance carries a unique reference designator; each net a unique name;
and the design database treats collisions as errors. This sounds trivial
and is not: unique, stable identity is what makes diff, review, archive,
and cross-reference possible across a design's lifetime, and its absence
is what makes informal diagrams unauditable. CAS PCBA 8's board structure
--- immutable, identity-bearing records for components, ports, nets, test
points, calibrations, and feedback declarations, with uniqueness enforced
by its checker --- is the netlist discipline, borrowed whole.

\subsection{Design-rule checking as an industrial discipline}\label{ch02:eda:drc}

Design-rule checking (DRC) arose in physical design: fabrication processes
publish geometric rules (minimum widths, spacings, enclosures), and a DRC
tool checks a layout against the rule deck exhaustively, emitting a
report of violations with locations. Over decades this practice hardened
into a sign-off culture with several transferable characteristics.

First, \emph{the rule deck is versioned and closed}: checks come from a
defined catalog with stable identifiers, so that ``DRC clean'' is a
statement relative to an exact deck, and two runs against the same deck
and the same design are comparable. Second, \emph{checking is exhaustive
and deterministic}: the tool checks every instance of every rule, and its
report is a function of its inputs, not of operator attention. Third,
\emph{clean is a gate, not a grade}: a design with violations does not
proceed, and waivers --- deliberate acceptances of specific flagged
violations --- are themselves recorded artifacts with identities and
justifications, never silent. Fourth, \emph{the report is the product}:
what sign-off consumes is not the designer's assurance but the tool's
report bound to the design version it checked.

Electrical rule checking (ERC) extended the same discipline from geometry
to connectivity semantics, and its standard check families are worth
enumerating because they map so directly onto the monograph's subject
matter: outputs shorted to outputs (multiple drivers on one net); inputs
left floating (undriven loads); ports connected across incompatible
signal families (a power-kind pin into a signal-kind pin); fan-out and
loading violations; and unconnected required pins. ERC checks are shallow
individually and powerful collectively: each encodes a semantic
expectation about connectivity that a drawing cannot enforce, and
industrial experience is that they catch, mechanically and early, a class
of defect that no amount of careful drawing eliminates.

CAS PCBA 8's design-rule checker is a deliberate transplant of this
discipline: a closed, versioned diagnostic catalog (\texttt{PCBA8-*});
deterministic single-report semantics; errors that fail the board;
warnings that never decay into an implicit pass; and net-class checks ---
unknown endpoint owner, invalid driver count, missing net load,
signal-kind mismatch, unit mismatch, multiply-connected ports, required
ports unconnected --- that any ERC engineer would recognize at a glance.
The transplant is organizational. That DRC/ERC discipline works for
electronics establishes that fail-closed cataloged checking is a sound
engineering practice; it establishes nothing about whether any particular
CAS rule expresses a true constraint on cognition, and the monograph
never argues otherwise.

\subsection{Consistency checking between representations}\label{ch02:eda:lvs}

A third pillar of EDA sign-off practice, alongside DRC and ERC, deserves
mention because it names a check \emph{shape} the monograph reuses:
layout-versus-schematic (LVS) checking, which extracts the connectivity
actually present in a physical layout and compares it, device by device
and net by net, against the schematic netlist that the layout purports to
implement. LVS institutionalized the insight that when the same design
exists in two representations, their agreement is not to be assumed but
computed --- and that the comparison must be canonical (both sides reduced
to a normal form before matching) for the verdict to be meaningful. The
general shape --- two representations of one design, a canonicalization,
a mechanical equivalence check, and a fail-closed verdict --- is exactly
the shape of the round-trip obligations the candidate imposes on any
future CCDL grammar (parse the corpus, serialize canonically, and match),
and of the audit's comparison between the source's published grammar and
the source's published boards. In LVS terms, the monograph's
grammar-coverage finding is a failed cross-representation consistency
check on the source canon itself, executed at the level of surface
syntax under a disclosed, bounded method.

Waiver practice completes the sign-off picture. Real designs sometimes
carry deliberate rule violations --- a spacing waived for a matched
analog pair, a floating pin documented as intentionally unused --- and
mature flows do not handle these by relaxing the rule or hand-editing
the report. They record a \emph{waiver}: an identified, justified,
reviewed artifact that acknowledges the specific violation instance and
travels with the design. The report stays true (the violation is real);
the disposition is explicit; and an auditor can enumerate every place
where human judgment overrode the deck. The candidate's analog is its
explicit exception discipline --- for instance, declared low-gain
exceptions in feedback analysis --- in which every departure from the
conservative default is itself a typed, identified declaration rather
than a silent tolerance.

\subsection{Fail-closed as a cultural achievement}\label{ch02:eda:failclosed}

One aspect of EDA practice deserves separate emphasis because it is a
culture, not an algorithm. Mature checking tools treat \emph{inability to
check} as failure: a rule the tool cannot evaluate on a given structure
produces an error or an explicit unresolved status, never a silent pass.
The alternative --- treating unevaluated as clean --- is seductive
(reports look better) and catastrophic (the report's silence becomes
uninterpretable). The fail-closed convention is what makes a clean report
mean something. CAS PCBA 8 adopts it globally: unsupported feedback
analysis yields an unresolved-safety diagnostic; unsupported simulation
behavior fails closed; an unresolved reachability bound is not treated as
reachable. The convention costs green checkmarks and buys interpretability,
and the monograph regards the trade as settled by several decades of
industrial precedent.

\section{Graph algorithms for circuit analysis}\label{ch02:graphs}

\subsection{Circuits as directed graphs}\label{ch02:graphs:directed}

A netlist induces graphs, and most connectivity analyses are graph
algorithms. For the analyses relevant to this monograph, the induced
object is a directed graph on components: an edge runs from component $u$
to component $v$ when some net is driven by a port of $u$ and loads a
port of $v$. Questions about signal flow become questions about this
graph: acyclicity determines whether a design can be evaluated in one
forward pass; cycles locate feedback; reachability delimits the influence
cone of an input; and topological order, when it exists, is an
evaluation schedule.

\subsection{Depth-first search and strongly connected components}\label{ch02:graphs:tarjan}

The classical toolkit here is depth-first search, and its crown result for
present purposes is Tarjan's algorithm \cite{tarjan}, which computes the
strongly connected components (SCCs) of a directed graph in a single
depth-first traversal, in time linear in the number of vertices and
edges. A strongly connected component is a maximal vertex set within
which every vertex reaches every other; contracting each SCC to a single
node yields the condensation of the graph, which is always acyclic. The
decomposition is exactly the right primitive for feedback analysis:
every directed cycle lies within one SCC, so the SCCs of size greater
than one (together with trivial SCCs bearing self-loops) enumerate the
regions of a design where feedback exists, and the acyclic condensation
orders everything else for straight-line evaluation.

The algorithm itself repays a paragraph of description, both because its
economy is instructive and because the monograph's checker relies on its
completeness. A depth-first search from each unvisited vertex assigns
every vertex a discovery index and maintains a stack of vertices whose
component is not yet determined. For each vertex $v$ the algorithm
computes a \emph{lowlink}: the smallest discovery index reachable from
$v$ by following tree edges and at most one back or cross edge into the
stack. When the search finishes a vertex whose lowlink equals its own
discovery index, that vertex is the root of a strongly connected
component, and the component is exactly the set of vertices above it on
the stack, which are popped together. Every vertex is pushed once and
popped once, and every edge is examined a constant number of times ---
hence linearity --- and the components are emitted in reverse
topological order of the condensation, so a single pass yields both the
feedback regions and an evaluation order for the acyclic remainder
\cite{tarjan}. That last property is quietly load-bearing for the
architecture: the same classical computation that finds every cycle also
schedules everything that is not a cycle, which is why a checker and a
simulator can share one graph analysis.

Two properties of this result matter methodologically. First, it is
\emph{complete}: linear-time SCC computation finds all feedback regions,
so a checker built on it can assert the absence of undeclared cycles, not
merely the presence of found ones. Second, it is \emph{cheap}: linearity
means the analysis scales to any design a database can hold, so there is
no engineering pressure to sample or skip. CAS PCBA 8's feedback check
class builds the directed component graph, computes SCCs by the classical
method \cite{tarjan}, and then applies its own conservative policy to
each realizable declared cycle: an exact feedback declaration, a declared
unit delay, an existing dimensionless calibration, and a sub-unity gain
bound, with explicit declared exceptions and with stateful or stochastic
loops unsupported and failing closed. The graph theory is classical; the
policy layered on it is the candidate's own contract, and its
conservatism is a design choice, not a theorem.

\subsection{Bound propagation on acyclic structure}\label{ch02:graphs:bounds}

The condensation's acyclicity also underwrites a second analysis pattern
used later in the monograph: conservative bound propagation. On an
acyclic evaluation order, an upper bound on each component's output can
be computed from upper bounds on its inputs and the component's declared
behavior (a clamp caps a bound; a monotone non-amplifying stage preserves
it; a gain stage scales it), and the composition of per-stage bounds
along a path yields a sound bound at the path's end. Interval-style
propagation of this kind is standard static-analysis practice; its
soundness depends on each per-stage bound being sound and on the
traversal respecting dependency order. The design-rule checker's
threshold-reachability class --- which flags a safety-critical threshold
that provably exceeds the best upstream bound --- is an application of
this standard pattern, and its verdict on the source's Cognitive Safety
Board is developed in later chapters.

\section{Simulation verification versus operational validation}\label{ch02:vandv}

\subsection{Two questions that must not be merged}\label{ch02:vandv:questions}

The modeling-and-simulation community maintains a foundational distinction,
surveyed authoritatively by Sargent \cite{sargent}, between
\emph{verification} and \emph{validation}. Verification asks: does the
implemented model correctly implement its conceptual specification --- is
the code a faithful realization of the model the developers intended?
Validation asks: is the conceptual model an adequate representation of the
problem entity for the model's intended purpose --- does the model,
however faithfully implemented, track the world well enough to use? The
two questions have different subjects (software versus world), different
evidence (testing, tracing, and code analysis versus comparison against
system data, expert scrutiny of behavior, and predictive trials), and
different failure modes (bugs versus wrong models), and Sargent's
framework further distinguishes conceptual-model validity, data validity,
and operational validity as separate obligations within the overall
modeling process, each requiring its own evidence at its own stage.

Sargent's survey is valuable to this monograph not only for the
distinction but for its catalog of \emph{validation techniques}, because
the catalog shows concretely how different the two evidence regimes are.
Verification draws on software-engineering technique: structured
walkthroughs, traces of program logic, testing against the conceptual
model's specification. Validation draws on a different toolbox entirely:
comparison of model output against data from the real system where such
data exist; comparison against other validated models; degenerate-case
and extreme-condition tests (does the model behave sensibly when inputs
are driven to zero or to their limits?); internal-validity checks on
stochastic replications; face validation by domain experts examining
behavior; historical data validation with held-out samples; predictive
validation against outcomes not used in construction; and animation or
operational graphics reviewed by people who know the modeled system.
Almost every item on the validation side presupposes access to the
modeled system, to its data, or to experts on it --- which is exactly
why validation evidence cannot be manufactured inside the development
loop, and why a project that possesses only development-loop evidence
must say so in so many words. Sargent's framework also emphasizes that
validity is always relative to an \emph{intended purpose} and required
accuracy: a model is not valid in the abstract but adequate-for-use, and
a validity claim without a stated purpose is not yet a claim.

The community codified this distinction because merging the questions
produces a specific, recurring error: a model that reproduces its own
specification perfectly is treated as if it had been shown to represent
reality. Verification evidence is comparatively cheap --- it can be
manufactured entirely inside the development loop --- while operational
validation requires access to the modeled system, experimental design,
and statistical care. Wherever the expensive evidence is scarce, the
cheap evidence expands to fill the vacuum unless the vocabulary itself
resists.

\subsection{Position of the present work}\label{ch02:vandv:position}

In Sargent's vocabulary \cite{sargent}, everything green in CAS PCBA 8 at
the bound source state is verification evidence: bounded regression
suites showing that the checker, simulator, and adapters implement their
declared contracts on their tested surfaces. Operational validation ---
evidence that any circuit-shaped cognitive model tracks any cognitive
system --- has not been attempted, is not claimed, and appears in the
monograph only as an open obligation in the research program. The
model-versus-world separation of Chapter~\ref{ch01:introduction} is the
verification/validation distinction enforced mechanically: by binding
every trace to a research boundary and typing every receipt as software
evidence, the architecture makes the cheap-evidence-expansion error a
detectable category violation rather than a rhetorical drift. The
distinction itself is the simulation community's, not this project's;
the project's contribution, if any, is only the mechanical enforcement.

\section{Property-based testing}\label{ch02:pbt}

\subsection{From examples to properties}\label{ch02:pbt:idea}

Conventional unit testing checks a program against hand-chosen
input/output examples. Its blind spot is structural: the examples test
the cases the author thought of, and defects live disproportionately in
the cases nobody thought of. Property-based testing, introduced to
practice by Claessen and Hughes's QuickCheck \cite{quickcheck}, inverts
the economy: the author writes a \emph{property} --- a predicate that
should hold for all inputs of a type --- and the framework generates
random inputs from that type's distribution, checks the property on each,
and reports the first counterexample. Two further ideas complete the
method. \emph{Generators} are first-class: the author controls the
distribution of random inputs, composing generators for structured data
so that random boards, random expressions, or random traces are
constructible. \emph{Shrinking} makes counterexamples usable: on
failure, the framework searches for a smaller input that still fails,
typically presenting a near-minimal witness rather than the sprawling
random structure that first triggered the fault.

Concreteness helps here, because the properties relevant to this
monograph are of a recognizable family. A round-trip property says that
serialization inverts parsing: for every generated board $b$,
$\mathrm{parse}(\mathrm{serialize}(b)) = b$ under canonical equality. A
determinism property says that a computation is a function of its
declared inputs: for every generated board, two runs of the checker
yield byte-identical reports. An invariance property says that an
equivalence the design intends is real: for every board and every
permutation of its component declaration order, the canonical digest is
unchanged. A metamorphic property relates outputs across related inputs
without knowing either output in advance: adding an undriven net to any
clean board must produce exactly one new missing-load diagnostic. Each
property is a few lines to state; each quantifies over a space no
example suite can cover; and each, under random generation with
shrinking, yields on failure a small witness board that a developer can
read. The QuickCheck lesson \cite{quickcheck} is that
specifications of this shape are cheap to write and disproportionately
effective to check --- and the corresponding honesty obligation is that
until such a campaign is run, invariants of this shape are contracts
with example-level evidence only.

The method's track record established a general lesson: for software
whose specification is expressible as invariants --- parsers that must
round-trip, checkers that must be deterministic, orders that must be
total, encodings that must invert --- randomized property checking finds
defects that example suites systematically miss, at low authoring cost.
The lesson comes with an honest converse: property testing samples the
input space and therefore establishes confidence, not proof; a passed
property is evidence proportional to the generator's coverage of the
space, and a biased generator can starve the very region where the
defect lives.

\subsection{Role in this work}\label{ch02:pbt:role}

CAS PCBA 8's contracts are invariant-shaped throughout --- one
deterministic DRC report per board; trace determinism under a declared
profile; digest stability under canonical serialization; admissibility
tables that must be closed under their own composition rules --- and are
therefore natural property-testing targets. The monograph is explicit
about status: at the bound source state, the architecture's evidence is
bounded example-style regression suites, and a property-based/random
generation campaign over DRC and simulation invariants is a scheduled
step of the research program \cite{quickcheck}, not an accomplished one.
The citation is to the method; no property-testing result is claimed,
and QuickCheck is not authority for the separately listed
mutation-testing obligation, which requires its own sources and
specification.

\section{Provenance}\label{ch02:prov}

\subsection{The provenance vocabulary}\label{ch02:prov:vocab}

When computational results feed decisions, the question ``where did this
number come from?'' needs a machine-readable answer. The W3C PROV family
of specifications, with PROV-O as its ontology realization \cite{prov},
standardizes a small core vocabulary for exactly this: \emph{entities}
(things with fixed state --- documents, datasets, results),
\emph{activities} (processes that occur over time and act on entities),
and \emph{agents} (things bearing responsibility), connected by typed
relations --- an entity \emph{was generated by} an activity, an activity
\emph{used} entities, an entity \emph{was derived from} another entity,
an entity or activity \emph{was attributed to} or \emph{associated with}
an agent. The apparent modesty of the vocabulary is its strength: with
generation, use, derivation, and attribution made explicit and typed,
lineage questions become graph queries --- what did this result depend
on, transitively? which results does this defective input contaminate?
--- and independent systems can exchange lineage without bilateral
format agreements.

The vocabulary scales in both directions. Downward, PROV supports
specialization and qualified relations, so that a bare ``was derived
from'' can be refined --- which activity performed the derivation, using
which entity, in which role --- when finer lineage matters. Upward, PROV
supports \emph{bundles}: provenance descriptions that are themselves
entities, so that one can record the provenance of a provenance record
--- who asserted this lineage, generated by what process, when. The
bundle idea matters for any architecture whose receipts might be
mistaken for authority: it makes ``this receipt was generated by the
same process it vouches for'' an expressible, queryable fact rather than
a buried circumstance, which is precisely the fact one needs surfaced to
refuse self-certification.

Two design points of PROV matter for the borrowing at hand. First,
derivations form a directed structure whose acyclicity within a record
set is a checkable well-formedness condition: nothing is its own
ancestor. Second, provenance assertions are claims by agents, not
oracles; PROV records who says a derivation held, and reconciling or
trusting such claims is explicitly out of the vocabulary's scope. The
vocabulary makes lineage sayable and queryable; it does not make it
true.

\subsection{Role in this work}\label{ch02:prov:role}

The canonical traces and measurement records of CAS PCBA 8 keep the same
separations the PROV vocabulary keeps --- results (entity-like) apart
from runs (activity-like) apart from derivations, with parent/child
measurement links validated for existence and acyclicity, and every
record bound to identities rather than descriptions. The monograph's
claim here is deliberately minimal and matches the source documents: the
record discipline \emph{resembles} the entity/activity/derivation
separation of W3C PROV \cite{prov}, and no PROV-O conformance is
claimed. The resemblance is cited to locate the design in a known
methodological tradition; conformance would be a checkable claim the
project has not earned and therefore does not make. The second PROV
design point --- provenance as claim, not oracle --- reappears in the
architecture's rule that locally asserted calibration metadata cannot
authenticate itself: recording lineage in valid shape never converts the
recorded content into evidence.

\section{Arithmetic control: floating point and decimal contexts}\label{ch02:arith}

\subsection{Why arithmetic is a reproducibility concern}\label{ch02:arith:why}

A deterministic, replayable computation contract must specify its
arithmetic, because ``the same computation'' is underdetermined without
it. IEEE Std 754 \cite{ieee754} is the standard that made binary
floating-point arithmetic specifiable: it defines the binary formats
(binary64 being the ubiquitous double precision), the rounding-direction
attributes, the required correctly rounded basic operations, exceptional
values and their propagation, and the status flags that record inexact,
overflow, underflow, and invalid events. The standard's achievement is
frequently misstated, and the accurate statement matters here: IEEE 754
makes each \emph{individual operation} exactly specified, but a
\emph{computation} composed of many operations remains sensitive to
evaluation order, compiler transformations, intermediate precisions,
library function variation, and parallel reduction order. Bitwise
reproducibility across toolchains is achievable but is an engineering
discipline built on the standard, not a free consequence of it. Floating
point is also, famously, non-associative in the small: rounding makes
$(a+b)+c$ and $a+(b+c)$ potentially different, which is exactly why
evaluation order is part of a computation's identity.

Decimal arithmetic under an explicit context object is the other
established control regime the monograph's subject draws on. In a
context-based decimal system, every operation executes under a context
that fixes precision, rounding mode, and exceptional-condition handling;
the context is data, so a computation can construct, version, and record
the context it ran under, and two parties holding the same context and
the same operation sequence obtain identical digit strings. The
trade-offs between the regimes are conventional: binary64 buys hardware
speed and a vast numerical literature at the price of representation
error for decimal quantities and order sensitivity; explicit decimal
contexts buy exact decimal representation and recordable, versioned
arithmetic identity at the price of speed.

\subsection{Role in this work}\label{ch02:arith:role}

CAS PCBA 8 uses both regimes, each disclosed. Candidate arithmetic runs
under one immutable, versioned decimal context created from primitive
constants for each operation, with ambient caller precision made
irrelevant and the context identity bound into reports and results ---
an application of the context-control pattern, with IEEE 754
\cite{ieee754} cited as the external reference point for explicit
arithmetic control while the candidate's decimal profile remains its own
software contract. Separately, the disclosed binary64 implementation of
the Fisher--Rao distance (Section~\ref{ch02:infogeo}) is characterized in
the source documents as a numerical pseudometric at its resolution ---
an honest floating-point disclosure of exactly the kind this section's
background predicts is necessary: near-identical profiles may be
numerically indistinguishable, so the implementation does not claim to
separate all distinct points, and independent numerical audit remains an
open obligation.

\section{Information geometry on probability simplices: a primer}\label{ch02:infogeo}

This section develops, at primer depth, the geometry that the
architecture's mixed-profile state addresses inhabit. The material is
classical; modern book-length treatments are Amari \cite{amari} and Ay et
al.\ \cite{ay}, and the founding papers are Fisher's on the information
quantity \cite{fisher} and Rao's on the induced distance \cite{rao}.

\subsection{The open probability simplex}\label{ch02:infogeo:simplex}

Fix $n$ mutually exclusive categories. A probability distribution over
them is a point of the closed simplex
$\Delta_{n-1}=\{p\in\mathbb{R}^n : p_k\ge 0,\ \sum_k p_k=1\}$; the
\emph{open} simplex $\Delta_{n-1}^{\circ}$ requires strict positivity,
$p_k>0$ for all $k$. The open simplex is an $(n-1)$-dimensional smooth
manifold, and restricting to it is not pedantry: the geometric structure
below degenerates on the boundary (where some $p_k=0$), and a state
space defined on the open simplex builds ``every category retains
nonzero mass'' into the type. A tangent vector $u$ at $p$ is a direction
of infinitesimal reweighting and satisfies $\sum_k u_k=0$, since motion
must preserve total mass.

\subsection{The Fisher--Rao metric}\label{ch02:infogeo:metric}

A Riemannian metric assigns each point an inner product on its tangent
space, thereby defining lengths of curves and distances between points.
The Euclidean metric inherited from $\mathbb{R}^n$ is available but
statistically arbitrary: it treats a mass shift of $0.01$ the same
whether it moves between categories holding $0.5$ each or between
categories holding $0.011$ and $0.989$, though the statistical
distinguishability of the two changes is utterly different. Fisher's
information quantity \cite{fisher} supplies the statistically motivated
alternative: for a parametric family, the Fisher information measures
the expected curvature of the log-likelihood --- how sharply data
discriminate nearby parameter values. Taking the categorical family with
the probabilities themselves as coordinates, the Fisher information
matrix induces the \emph{Fisher--Rao metric} on
$\Delta_{n-1}^{\circ}$:
\[
g^{\mathrm{FR}}_p(u,v)\;=\;\sum_{k=1}^{n}\frac{u_k\,v_k}{p_k},
\]
for tangent vectors $u,v$ at $p$. The $1/p_k$ weighting encodes the
statistical intuition exactly: reweighting a rare category is a longer
move, per unit of mass, than reweighting a common one. Rao's 1945 paper
\cite{rao} was the first to treat this induced Riemannian structure as a
geometry on the space of distributions and to propose its geodesic
distance as a measure of dissimilarity between populations --- the
founding act of the field now called information geometry
\cite{amari,ay}.

A two-category example makes the weighting tangible. On
$\Delta_1^{\circ}$, parameterize points by $p=(t,1-t)$ with $0<t<1$; a
tangent vector is $u=(s,-s)$. The Fisher--Rao squared length of $u$ at
$p$ is
\[
g^{\mathrm{FR}}_p(u,u)=\frac{s^2}{t}+\frac{s^2}{1-t}
=\frac{s^2}{t(1-t)},
\]
which is minimized at $t=\tfrac12$ and diverges as $t\to 0$ or
$t\to 1$: the same nominal mass shift $s$ is geometrically short between
balanced categories and arbitrarily long near the boundary. The
divergence at the boundary is the concrete face of the earlier remark
that the geometry degenerates when any $p_k=0$, and it is the
mathematical reason a state space built on this metric is typed over the
\emph{open} simplex: the metric itself refuses the boundary.

\subsection{The sphere map and the closed-form geodesic distance}\label{ch02:infogeo:geodesic}

The categorical Fisher--Rao geometry admits a classical closed form via
the square-root embedding. Map $p\mapsto \xi(p) =
2(\sqrt{p_1},\dots,\sqrt{p_n})$. Since $\sum_k(\sqrt{p_k})^2=1$, the
image $\xi(p)$ lies on the sphere of radius $2$ in $\mathbb{R}^n$, in
its strictly positive orthant; and a short computation shows the map is
an isometry from $(\Delta_{n-1}^{\circ}, g^{\mathrm{FR}})$ onto its
image with the sphere's round metric. The computation is one line of
calculus: since $\xi_k = 2\sqrt{p_k}$, an infinitesimal displacement
satisfies $d\xi_k = dp_k/\sqrt{p_k}$, so the Euclidean squared line
element of the image is
$\sum_k d\xi_k^2 = \sum_k dp_k^2/p_k$ --- which is exactly the
Fisher--Rao squared line element on the simplex. The curved statistical
geometry is, up to this change of variables, a piece of a round sphere. Geodesic distance on a sphere of
radius $r$ is $r$ times the angle between unit vectors, so the
Fisher--Rao geodesic distance pulls back to
\[
d(p,q)\;=\;2\arccos\Bigl(\sum_{k=1}^{n}\sqrt{p_k\,q_k}\Bigr).
\]
The argument of the arccosine, $\sum_k\sqrt{p_k q_k}$, is the
Bhattacharyya coefficient of $p$ and $q$; it equals $1$ exactly when
$p=q$ (giving distance $0$) and decreases as the distributions
separate. Because the positive orthant of the sphere has angular
diameter $\pi/2$ between its axis points, distances are bounded:
$0\le d(p,q)\le \pi$ on the closed simplex, with the maximum attained
only between distributions of disjoint support --- which the open
simplex excludes, so on $\Delta_{n-1}^{\circ}$ the distance is always
strictly less than $\pi$. The formula makes the earlier metric intuition
quantitative and gives implementations a two-line evaluation: an inner
product of square roots and an arccosine.

Numerically, the formula has a disclosed sharp edge that
Section~\ref{ch02:arith} predicts: for very close $p$ and $q$ the
Bhattacharyya sum approaches $1$, where the arccosine's derivative
diverges, so binary64 evaluation loses relative accuracy exactly in the
near-coincident regime. A binary64 implementation is therefore honestly
described as a numerical pseudometric at its resolution --- it may fail
to separate distinct but extremely close profiles --- and that is
precisely the characterization the architecture's source documents
adopt.

\subsection{Invariance, and what the geometry does and does not say}\label{ch02:infogeo:invariance}

The Fisher--Rao metric is not merely one defensible choice among many;
classical results in information geometry (associated with Chentsov and
treated in the modern monographs \cite{amari,ay}) single it out, up to
scale, as the Riemannian metric on these spaces invariant under the
natural statistical maps between categorical models. Invariance is the
principled reason a system committed to representing distributional
states might adopt this metric rather than the Euclidean one: the
geometry then reflects statistical structure rather than an arbitrary
coordinate convention.

The boundary statement for this monograph is the same as every other
section's. That the Fisher--Rao geometry is well-founded mathematics
\cite{fisher,rao,amari,ay} makes it a sound \emph{setting} for
strictly positive normalized profiles over a finite cell set --- which
is how CAS PCBA 8 uses it, as the named metric profile for its
40-dimensional mixed-state addresses. The mathematics does not identify
a learned model's internal state with a CAS state, does not make a
profile over ND-cells an empirical measurement of anything, and does
not convert geodesic distance between two software-constructed profiles
into a cognitive quantity. The geometry constrains the model's internal
bookkeeping; it asserts nothing about the world.

\section{A note on causal claims}\label{ch02:causal}

One further piece of background belongs in this survey because its
\emph{absence} from the architecture's claims is load-bearing. The
modern mathematics of causal inference, systematized by Pearl
\cite{pearl}, draws a sharp line between associational quantities
(computable from observed joint distributions) and causal quantities
(defined by interventions, requiring a causal model and identification
assumptions to connect to data). A directed graph whose edges depict
influence is, by itself, on the associational side of the line: drawing
an arrow, propagating a bound along it, or simulating a transfer rule
across it does not establish that intervening on the source would move
the target. CAS PCBA 8's boards are directed graphs with declared
transfer semantics, and the temptation to read them causally is
therefore structural. The architecture's position, developed with the
threats to validity later in the monograph, is Pearl's line enforced as
a boundary: graph connectivity, gain composition, prediction, and
residual localization do not establish causation; causal claims would
require explicit causal models, identification assumptions, and
interventions \cite{pearl}; and none is admitted.

\section{What was deliberately not borrowed}\label{ch02:notborrowed}

A survey of borrowings is incomplete without its negative space: several
prominent instruments of the surveyed fields were considered by the
architecture's design lineage and \emph{not} adopted, and the reasons are
themselves methodological content.

From hardware practice, the architecture does not borrow \emph{physical
modeling}. Electronic simulation's deepest layer --- device physics,
parasitic extraction, timing derived from geometry and process ---
presupposes that the modeled quantities are physical quantities with
instruments that measure them. Importing any of it would contradict
Boundary~\ref{ch02:bnd:background}'s companion discipline in
Chapter~\ref{ch01:introduction}: there is no physical profile in the
candidate, so there is nothing for device physics to be about. The
borrowing stops, deliberately, at the design-database layer ---
description, connectivity, rules, traces --- which is the layer that is
physics-free.

From formal methods, the architecture does not (yet) borrow \emph{proof}.
Model checking, theorem proving, and verified compilation all offer
stronger guarantees than cataloged checking and bounded suites, and
nothing in the architecture's design forecloses them. They are absent
because the honest sequencing runs the other way: proof obligations are
worth their cost when the specification is stable and complete, and the
monograph's own source audit shows the specification layer (the CCDL
grammar) is neither. Proving theorems about a grammar that does not
cover its corpus would harden the wrong artifact. The research program
places grammar completion and round-trip conformance before any formal
escalation, and until then the candidate claims exactly what bounded
suites support.

From statistics and machine learning, the architecture does not borrow
\emph{fitting}. Nothing in the candidate estimates parameters from data,
and this is a boundary, not a gap: a fitted parameter is a measurement
claim in disguise, and the design-value-versus-measurement separation
requires that the measurement ledger stay empty until an authenticated,
preregistered pathway fills it. The information geometry of
Section~\ref{ch02:infogeo} is adopted as a metric structure for typed
addresses, not as an estimation apparatus.

Each refusal could be reversed by future work with the appropriate
evidence discipline --- a physical profile with instruments, a proof
layer over a completed grammar, a preregistered estimation pathway ---
and the research program names the preconditions. The refusals are
recorded here so that their reversal, if it comes, is a visible event
with a visible warrant rather than a drift.

\section{Summary: the borrowings and their limits}\label{ch02:summary}

Table~\ref{ch02:tab:map} condenses the chapter. Each row names a
background item, the methodological lesson its home community
established, the CAS PCBA 8 mechanism that borrows the lesson, and ---
the column this monograph exists to keep honest --- what the borrowing
does not license.

\begin{longtable}{@{}p{2.6cm}p{3.4cm}p{3.4cm}p{4.0cm}@{}}
\caption{Background borrowings and their limits.}\label{ch02:tab:map}\\
\toprule
Background & Established lesson & Borrowing mechanism & Not licensed \\
\midrule
\endfirsthead
\toprule
Background & Established lesson & Borrowing mechanism & Not licensed \\
\midrule
\endhead
\bottomrule
\endfoot
VHDL \cite{vhdl} & Formal machine-readable design text with typed
interfaces and defined execution semantics & Typed board model; closed
catalogs; deterministic simulator & Any validity of CAS content written
in the language \\
Statecharts \cite{harel} & Behavioral diagrams need explicit state,
hierarchy, concurrency, communication semantics & Declared statefulness
and time semantics; declared feedback; fail-closed on undeclared
dynamics & Any cognitive interpretation of the states \\
EDA practice (general) & Netlist identity discipline; versioned closed
rule decks; fail-closed sign-off & Board identity records;
\texttt{PCBA8-*} catalog; errors fail the board & Truth of any specific
CAS rule about cognition \\
Tarjan \cite{tarjan} & Linear-time complete SCC decomposition & Feedback
region enumeration; acyclic evaluation order; bound propagation &
Soundness of the candidate's feedback policy beyond its stated
conservatism \\
Sargent \cite{sargent} & Verification and operational validation are
distinct obligations & Evidence taxonomy; trace-bound research boundary
& Any operational validation of any board \\
QuickCheck \cite{quickcheck} & Randomized property checking finds
defects example suites miss & Scheduled research-program campaign over
DRC/simulation invariants & Any completed property-testing result;
mutation-testing authority \\
PROV \cite{prov} & Typed generation/use/derivation vocabulary makes
lineage queryable & Trace and measurement record discipline; derivation
acyclicity checks & PROV-O conformance; truth of recorded content \\
IEEE 754 \cite{ieee754} & Arithmetic must be explicitly specified for
reproducibility & One immutable decimal context; disclosed binary64
pseudometric caveat & Bitwise reproducibility beyond the declared
contract \\
Fisher, Rao, Amari, Ay \cite{fisher,rao,amari,ay} & Invariant Riemannian
geometry on the open simplex with closed-form geodesic distance & Named
Fisher--Rao profile for 40-cell mixed addresses & Identification of
model states with CAS states; empirical meaning of distances \\
Pearl \cite{pearl} & Association and intervention are distinct; causal
claims need models and assumptions & Causal-claim quarantine in the
threat model & Any causal reading of board structure or simulation \\
\end{longtable}

The chapter closes where it began. Everything above is the
accomplishment of other communities, surveyed because the architecture
of this monograph is assembled from their methods. The assembly is the
subject of the chapters that follow; its evidence begins with the
versioned source basis fixed in Chapter~\ref{ch03:source}; and at no
point does the strength of the surveyed foundations transfer, by
association, to the cognitive content modeled on top of them.
